The oscillation occurs due to the fact that the center of mass (CM) of the
Solar System does not lie
% sic: source has "lieexactly"
exactly at the center of the Sun, but in another position around which the
center of the Sun revolves.

% note: the source uses the astronomical symbol for Jupiter as subscript;
% rendered here as "Jup".
Given $a_{\mathrm{Jup}} = 5.204$~AU, the distance from CM to the center of the Sun
is:
\[ x_{CM} = \frac{\frac{1}{1047}}{1+\frac{1}{1047}}\,5.204 = 0.004966\ AU, \]
\hfill(3 pt: position of CM)

which is the amplitude of the movement.

If the distance to Barnard's Star is $d_{bar} = 1.834$~pc, we can write, for
the angular amplitude:
\[ \alpha \approx \frac{x_{cm}}{d_{bar}} = 0{,}002708'' \]
\hfill(3 pt: angular amplitude)

As for the period, it is the period of Jupiter's orbit:
\[ \frac{T_{\mathrm{Jup}}^2}{a_{\mathrm{Jup}}^3} = \frac{T_\oplus^2}{a_\oplus^3}
   = 1\,\frac{year^2}{AU^3}
   \quad\therefore\quad T_{jup} = 11{,}87\ years \]
\hfill(4 pt: period of oscillation)

Obs: the difference of Jupiter's mass in the Kepler's Third Law represents a
difference in less than 0.01 year at the period of Jupiter.
